% !TEX root = ../PhD Thesis.tex
\chapter{The SG ``Indirect" Proteome}
\label{chap:ChapterII}

Unlike \gls{rna}, proteins display a more uniform sedimentation behaviour during \gls{sg} isolation, with their enrichment being less strongly influenced by stress-dependent redistribution. This likely reflects a fundamental difference in how RNA and proteins participate in stress-induced condensates: mRNAs are broadly remodelled into ribonucleoprotein assemblies during stress, leading to coordinated changes in their partitioning, whereas protein recruitment is more selective and restricted to specific subsets of factors associated with these complexes\cite{glauninger_transcriptome-wide_2025}. This is reflected in the high consistency observed across proteomic studies\cite{nunes_msgp_2019, millar_new_2023}. Notably, two independent compiled databases of \gls{sg}-associated proteins report highly overlapping top candidates and converge on the conclusion that approximately 50\% of them are \glspl{rbp}\cite{nunes_msgp_2019, millar_new_2023}.
\\
\\
Yet, the presence of a protein in \gls{sg}s does not imply that all of its cellular copies are sequestered within \glspl{sg}. In fact, as discussed previously, even canonical \gls{sg} markers such as \gls{g3bp1} show only partial localisation, with estimates suggesting that approximately 20\% of the total protein pool is associated with \glspl{sg}\cite{somasekharan_g3bp1-linked_2020}. The functional consequences of this partial sequestration remain unclear, as \glspl{rbp} are highly multifunctional and participate in numerous aspects of \gls{rna} metabolism (see section 1.2.2).
\\
\\
Importantly, this redistribution of \glspl{rbp} implies a reduction in their effective availability in other cellular compartments, rather than a simple gain of function within \glspl{sg}. In particular, sequestration of \glspl{rbp} may alter their nuclear-cytoplasmic balance, potentially reducing their participation in nuclear processes such as pre-mRNA splicing. Given the central role of many \glspl{rbp} in splice site recognition and spliceosome assembly, such changes could lead to measurable alterations in splicing patterns, providing a rationale for investigating ``SG-induced" differential splicing.

\section{Materials and Methods}

\subsection{Dataset Procurement and List}

For dataset procurement, we queried the \gls{geo} Datasets repository using the term ``stress granule", filtering for expression profiling by high-throughput sequencing and restricting results to human samples, following a strategy similar to that described in Section 3.1.1. In this case, however, we specifically targeted studies including both stressed and unstressed conditions, as well as experimental setups with and without \gls{sg} formation.
\\
\\
Because the objective was to perform differential splicing analysis, only datasets with sufficient sequencing depth were considered (At least 5 \glspl{gb} of bases sequenced per sample, in all samples). High read depth is particularly critical for splicing analyses, as exon–exon junction reads and low-abundance isoforms are often underrepresented, making shallow sequencing insufficient to reliably capture alternative splicing events. Based on these criteria, we selected the dataset GSE173953\cite{paget_stress_2023}, which includes \gls{u2os} and A549 \gls{g3bp1} \gls{wt} and \gls{ko} cells exposed to \gls{dsrna} as a proxy for viral infection. Importantly, \gls{g3bp1} knockout cells do not form \glspl{sg} under these conditions, enabling direct comparison between stressed cells with and without SG formation. This design allows the isolation of \gls{sg}-specific effects from general stress responses. Datasets GSE138988\cite{matheny_transcriptome-wide_2019}, and GSE119977\cite{matheny_rna_2021, matheny_corrigendum_2023} had the necessary sample types, but lower read coverage (Less than 5 \glspl{gb} sequenced per sample, in all samples) than GSE173953\cite{paget_stress_2023} and were thus discarded.

\subsection{Correlations}

Unless otherwise indicated, correlations shown in scatter plots and heatmaps were calculated using the t-statistic from differential expression, with Spearman's correlation coefficient $\rho$ and the associated test's p-value.

\subsection{Differential Splicing Analysis}

Raw sequencing reads (FASTQ format) were first trimmed using Trimmomatic (v0.39)\cite{bolger_trimmomatic_2014} in \gls{pe} mode, and subsequently processed with VAST-TOOLS (v2.5.1)\cite{tapial_atlas_2017}. Alignment was performed against the human reference genome \gls{hg38} using Bowtie (v1.2.2)\cite{langmead_ultrafast_2009}, following the standard VAST-TOOLS pipeline\cite{tapial_atlas_2017}.
\\
\\
After running the \texttt{vast-tools align} and \texttt{vast-tools combine} steps, differential splicing analysis was conducted using the betAS R package (v1.2.1)\cite{ascensao-ferreira_betas_2024}, with modifications to the default filtering procedure. By default, the \texttt{alternativeEvents} function excludes events in which at least one sample does not meet the \gls{psi} thresholds (1 $<=$ \gls{psi} $<=$ 99). However, applying these thresholds resulted in the inclusion of a large number of events that appeared constitutive, showing minimal variation rather than true alternative splicing.
\\
\\
To improve specificity, we incorporated prior information from the VASTDB MAIN PSI table\cite{tapial_atlas_2017}. For exon-skipping events, we computed the average \gls{psi} and variance across tissues and cell types. Events were classified as alternative if they satisfied at least one of the following empirical criteria: variance $>=$ 750, or mean \gls{psi} between 12.5 and 87.5. Only events meeting these criteria were retained for downstream filtering of the VAST-TOOLS output.
\\
\\
For \gls{shrna} experiments, data from the ENCODE project (HepG2 and K562 cells treated with scrambled \gls{shrna} controls)\cite{the_encode_project_consortium_integrated_2012} were processed jointly using this modified betAS pipeline. The same approach was applied to the GSE173953 dataset, with \gls{u2os} and A549 samples analysed together. Although this joint processing strategy reduced the total number of detected differential splicing events, it increased the robustness and reproducibility of the identified events.
\\
\\
A more in depth description of the integration of \gls{eclip} data with these differential splicing events is provided in a later chapter.

\subsection{\textit{In Silico} Resources}

All analyses were performed on the workstation \textit{lobito} at the \gls{gimm} data center, which is equipped with 19 \gls{tb} of \gls{ssd} storage, 768 \gls{gb} of RAM, 72 \gls{cpu} threads (\gls{intel} Xeon
Gold 6254 @ 3.10 \gls{ghz}), and an integrated Matrox G200eW3 \gls{gpu}. \gls{bash} was used to run vast-tools\cite{tapial_atlas_2017}, FASTQC\cite{andrews_fastqc_2010}, MultiQC\cite{ewels_multiqc_2016}, STAR\cite{dobin_star_2013}, and featureCounts\cite{liao_featurecounts_2014}. Downstream analyses were conducted in R (v4.1.2) using RStudio Server (v1.4.1717). The eCLIPSE tool (used for the \gls{eclip} splicing analyses, will be introduced in more detail next chapter) includes steps that were conducted in Python 3 (v3.10.11), and in R (v4.1.2). 


\subsection{U2-OS Growth Conditions}

Human osteosarcoma \gls{u2os} cells were cultured in DMEM (high glucose) supplemented with 10\% fetal bovine serum. Cells were plated on glass coverslips in 24-well plates and used for all experimental procedures.
\\
\\
Reagents and material references are supplied in section 10.1.1.

\subsection{Stress conditions}

To examine mRNA localisation under stress, when 80\% confluency was reached, cells were exposed to 0.4 \gls{molar} sorbitol for 3 hours.
\\
\\
Reagents and material references are supplied in section 10.1.1.

\subsection{Immunofluorescence}

Following stress treatments, cells were washed three times with \gls{pbs} +/+ (\gls{pbs} supplemented with calcium and magnesium). Fixation was performed using 4\% formaldehyde in \gls{pbs} +/+ (\gls{ph} 7.4) for 20 minutes at room temperature, followed by three additional washes with \gls{pbs} +/+. Cells were then permeabilised with 0.1\% Triton X-100 in \gls{pbs} +/+ for 10 minutes at room temperature and washed three times with unsupplemented \gls{pbs}.
\\
\\
Glass coverslips were removed from 24-well plates and transferred to a humid chamber containing 20 \gls{ul} droplets of blocking buffer (5\% goat serum in \gls{pbs}). Coverslips were inverted onto the droplets (cell side facing down) and incubated for 60 minutes at room temperature.
\\
\\
Primary antibodies against TIA1 were diluted 1:200 in blocking buffer and applied using the drop method (10 \gls{ul} per coverslip). Coverslips were incubated for 2 hours at room temperature in the humid chamber, followed by three washes with unsupplemented \gls{pbs}.
\\
\\
Secondary antibodies were diluted 1:500 in blocking buffer together with \gls{dapi} (1:1000) and applied (10 \gls{ul} per coverslip) for 1 hour at room temperature in the dark. Coverslips were then washed three times in unsupplemented \gls{pbs} with agitation (5 minutes per wash).
\\
\\
For mounting, 2 \gls{ul} of Vectashield mounting medium was placed on each microscope slide, and coverslips were positioned cell-side down onto the drops. Coverslips were sealed with nail polish prior to imaging.
\\
\\
Reagents and material references are supplied in section 10.1.1.

\subsection{Microscopy and Image Analysis}

After fixation, cells were imaged on a ZEISS \gls{lsm} 880 confocal microscope, using a 63x/1.4 Plan-Apochromat \gls{dic} oil immersion objective. For each condition, a z-stack of a single field of view (1024x1024\gls{px}, with a \gls{px} size of 0.13 \gls{um}) was acquired, comprising 10–20 Z-planes, with a 0.48 \gls{um} step size, according to Nyquist, with a scan zoom of 1.0. TIA1 was excited using a 633 \gls{nm} laser, and \gls{dapi} with a 405 \gls{nm} laser. Images were observed in \gls{fiji}, with maximum intensity thresholds defined for each channel using sorbitol-stressed samples. These same parameters were then applied to all other conditions to ensure consistency across images. 

\section{Results and Discussion}

\subsection{Sequestration of TIA1 in SGs Recapitulates Splicing Alterations from TIA1 Depletion}

Since a large number of transcriptomic datasets had been compiled for \glspl{sg}, we used alternative splicing as a proxy to assess the functional consequences of \gls{rbp} sequestration into \glspl{sg}. Changes in splicing are detectable at the transcriptomic level and provide a measurable readout of \gls{rbp} activity. If \gls{sg} sequestration affects splicing regulation, it is likely that other \gls{rbp}-dependent processes are also impacted (Figure 4.1).
\\
\begin{figure}[H]
	\centering
	\includegraphics[width=1\linewidth]{"../../../Desktop/Paper Images/TESE_Fig5.1"}
	\caption[Rationale behind RBP function alterations driven by SG sequestration]{\textbf{Rationale behind RBP function alterations driven by SG sequestration.} (Left) In the absence of \glspl{sg}, \glspl{rbp} remain available in the nucleus (or are recruited to) to regulate splicing as needed. (Right) When \glspl{sg} form, \glspl{rbp} are sequestered within them, reducing their nuclear availability and impairing splicing regulation.}
	\label{fig:Figure5.1}
\end{figure}
\noindent
Direct analysis of individual splicing events, however, is constrained by transcript expression\cite{van_nostrand_large-scale_2020}. Alternative splicing can only be detected in expressed genes, and expression profiles vary substantially across cell types and conditions, limiting the identification of consistent trends\cite{van_nostrand_large-scale_2020}. To address this, we adopted a complementary approach based on \gls{rbp} binding profiles. By grouping exon-skipping events into categories based on their response to \gls{sg} assembly (increased inclusion, no significant change, or increased exclusion), we constructed RNA splicing maps to assess whether specific \glspl{rbp} exhibit position-dependent binding patterns associated with these regulatory outcomes. In these maps, the distribution of RBP binding (or proximity) relative to alternative exons is compared across each regulatory class, allowing identification of characteristic binding signatures linked to exon inclusion or skipping (an example is provided in Figure 1.2, and this process is explained more in depth in Chapter 9).
\\
\\
Therefore, we compiled a set of \glspl{rbp} known to localise to \glspl{sg}\cite{nunes_msgp_2019, millar_new_2023} and mapped their \gls{rna} targets using \gls{eclip} data. We then performed differential alternative splicing analysis between \gls{dsrna}-stressed and unstressed cells. This analysis identified seven \glspl{rbp} (\gls{tia1}, \gls{qki}, \gls{ddx3x}, \gls{rbm15}, \gls{khsrp}, \gls{pcbp1}, and \gls{ybx3}) with significant enrichment or depletion patterns.
\\
\\
To minimise confounding effects from expression-driven splicing changes, we excluded \glspl{rbp} whose expression differed between stressed and unstressed conditions in both \gls{sg}-competent and \gls{sg}-deficient cells\footnote{More copies of a protein, are naturally expected to lead to, usually, increased function. It is perfectly acceptable that SG assembly would lead to expression changes that would naturally impact splicing regulation. Nevertheless, our aim was to uncover if SG sequestration could directly impact RBP function, and as such, expression driven splicing alterations were excluded. Obviously, other types of post-transcriptional could occur, potentially changing splicing such as RNA modifications, but as we had no way to infer them, we could not control for them.} (Figure 4.2A). This filtering yielded five candidate \glspl{rbp} with binding alterations compatible with lower protein availability: \gls{khsrp}, \gls{qki}, \gls{tia1}, \gls{pcbp1}, and \gls{ybx3}.

\begin{figure}[H]
	\centering
	\includegraphics[width=1\linewidth]{"../../../Desktop/Paper Images/Figure9"}
	\caption[Alternative splicing alterations in SG competent and incompetent cells]{}
	\label{fig:Figure5.2}
\end{figure}

\clearpage


\begin{figure}[t]
	\captionsetup{labelformat=adja-page}
	\ContinuedFloat
	\caption[]{\textbf{Alternative splicing alterations in SG competent and incompetent cells.} (A) Volcano plot of \gls{dgea} between stressed and unstressed \glspl{wc} (X-Axis -  Log2 Fold-Change) and corresponding significance (Y-Axis – B-Statistic) in SG competent (Left) and incompetent (Right) U2-OS and A549 cell lines subjected to dsRNA as a stressor\cite{paget_stress_2023}; genes as circles, with RBPs with differential binding across event types labelled and in red. Vertical dashed lines correspond to -0.5 and 0.5 Log2\gls{fc}. Horizontal dashed line corresponds to B-statistic $=$ 0. Only genes with -2 $<$ Log2\gls{fc} $<$ 2 are plotted. (B-F) RNA binding maps of alternative exon splicing events in shRNA treated cells for KHSRP (B), PCBP1 (C), QKI (D), TIA1 (E), and YBX3 (F) obtained from ENCODE\cite{the_encode_project_consortium_integrated_2012} (Left), in SG competent cells (Middle) and in SG incompetent cells (Right). (Top) Normalised binding of RBP of interest to metagenomic regulatory sequences of alternative exons. The plot discriminates excluded (blue), included (red) or maintained (black) alternative exons upon RBP knockdown (Left) or upon dsRNA stress stimuli, in SG competent (Middle) and incompetent (Right) cells. The total number of events in each category can be found beneath the plot in brackets, preceded by the number of events for which there is at least one binding site for that RBP. (Bottom) False discovery rate (FDR) of excluded (blue) or included (red) normalised binding, in relation to maintained events. Vertical lines delimitate the following metagenomic splicing regulatory sequences, in order: last 50 nt of upstream constitutive exon; first 200 nt of upstream intron; last 200 nt of upstream intron; first 50 nt of alternative exon; last 50 nt of alternative exon; first 200 nt of downstream intron; last 200 nt of downstream intron; first 50 nt of downstream constitutive exon.}
\end{figure}
\noindent
To determine whether the observed splicing patterns were consistent with loss of \gls{rbp} function, we compared these profiles with splicing changes following \gls{shrna}-mediated \gls{kd} of each \gls{rbp} (Figure 4.2B-F - Left). \glspl{rbp} whose binding profiles resembled the \gls{kd}-induced splicing patterns were considered likely to be functionally affected by \gls{sg} sequestration (Figure 4.2B-F - Middle). Of the five candidates, only \gls{tia1} and \gls{ybx3} showed visually similar FDR distributions, which were further assessed using Kolmogorov–Smirnov test (p-value = 0.16, and p-value = 0.9, respectively).
\\
However, these similarities could reflect general stress-induced changes rather than \gls{sg}-specific effects. To address this, we repeated the analysis in \gls{g3bp1}/\gls{g3bp2} \gls{ko} cells exposed to \gls{dsrna}, a condition in which \glspl{sg} do not form\cite{paget_stress_2023} (Figure 4.2B-F - Right). Under these conditions, the similarity observed for \gls{tia1} was lost (Kolmogorov-Smirnov test with p-value $=$ 0.02), whereas \gls{ybx3} remained unchanged (Kolmogorov-Smirnov test with p-value $=$ 0.46). This suggests that \gls{tia1} function may be specifically affected by \gls{sg} assembly, whereas \gls{ybx3} is more likely influenced by stress independently of \gls{sg} formation.
\\
\\
To further investigate \gls{tia1} behaviour, we performed immunofluorescence under control and hyperosmotic stress conditions. Under basal conditions, \gls{tia1} is typically nuclear\cite{noauthor_human_nodate, thul_subcellular_2017}, but in our experiments it was predominantly cytoplasmic, likely reflecting baseline stress associated with cell culture (Figure 4.3A). No clear differences in overall localisation were observed between conditions, although \gls{tia1} consistently marked \glspl{sg}.
\\
\\
Beyond its role in \gls{sg} assembly\cite{gilks_stress_2004}, \gls{tia1} is a well-characterised splicing regulator\cite{del_gatto-konczak_rna-binding_2000} and has been implicated in neurological disorders\cite{mackenzie_tia1_2017}, as well as in viral infection and replication\cite{emara_interaction_2007}. Several viruses interact with \gls{tia1} and its paralogue \gls{tiar}, inhibiting \gls{sg} formation\cite{emara_interaction_2007}. Whether this represents a direct strategy to disrupt \gls{sg} function or an indirect consequence of modulating \gls{tia1} activity remains unclear (see section 1.1.5).
\\
\\
Importantly, gene expression (Figure 4.3B) and splicing (Figure 4.3C) changes following \gls{tia1} knockdown\cite{the_encode_project_consortium_integrated_2012} did not recapitulate those observed during \gls{sg} assembly in \gls{dsrna}-treated cells\cite{paget_stress_2023}. This indicates that \gls{tia1} sequestration alone is unlikely to fully explain the broader transcriptomic effects associated with \gls{sg} formation.
 

\begin{figure}[H]
	\centering
	\includegraphics[width=1\linewidth]{"../../../Desktop/Paper Images/TESE_Fig5.3"}
	\caption[Alterations of TIA1 function upon KD and SG assembly]{}
	\label{fig:Figure5.3}
\end{figure}

\clearpage


\begin{figure}[t]
	\captionsetup{labelformat=adja-page}
	\ContinuedFloat
	\caption[]{\textbf{Alterations of TIA1 function upon KD and SG assembly.} (A) Immunofluorescence images of U2-OS cells untreated or treated with sorbitol (0.4M). TIA1 is shown in white and nuclei (DAPI) in blue. Yellow arrows highlight representative SGs. Scale bar is shown on bottom right corner of each image, in yellow. (B) Scatter plot of t-statistics of DGEA in HepG2 and K562 cells treated with shRNA against TIA1 versus untargeted shRNA\cite{the_encode_project_consortium_integrated_2012} (X-Axis) and between U2-OS and A549 WT Cells (SG-Competent) exposed to dsRNA as a stressor\cite{paget_stress_2023} (Y-Axis); genes as circles; density contour lines in dark red, Spearman's correlation coefficient $\rho$ and associated test's p-value on top left corner. (C) Scatter plot of \gls{dpsi} of differential splicing analysis in HepG2 and K562 cells treated with shRNA against TIA1 versus untargeted shRNA\cite{the_encode_project_consortium_integrated_2012} (X-Axis) and between U2-OS and A549 WT Cells (SG-Competent) exposed to dsRNA as a stressor\cite{paget_stress_2023} (Y-Axis); genes as circles; density contour lines in dark red, Spearman's correlation coefficient $\rho$ and associated test's p-value on top left corner.}
\end{figure}
\noindent